% GENERATED BY src/python/lmi/tables/ext_asyl_mediators.py
% -- DO NOT EDIT
%==============================================================================%
% ext_asyl_mediators.tex
%
% Purpose : Erweiterung, Tabelle E2 -- Kursangebot x Asylstatus auf die fuenf
%           Mechanismen. Kern der Erweiterung: getestet wird nicht der Ertrag
%           der Sprache, sondern der Zugang zur Behandlung selbst.
% Source  : SOEPremote jobs 244836, 244837, 244840,
%           src/stata/remote/ext_asyl_mediators_primary.do,
%           ext_asyl_mediators_secondary.do, ext_asyl_participation.do
%           Rohlistings: results/transcripts/src_stata_remote_ext_asyl_*.eml
%           Transkript: results/comparisons/ext_asyl_mediators.md
% Needs   : booktabs, tabularx (beide bereits in bachelorarbeit.sty geladen)
% Usage   : % GENERATED BY src/python/lmi/tables/ext_asyl_mediators.py
% -- DO NOT EDIT
%==============================================================================%
% ext_asyl_mediators.tex
%
% Purpose : Erweiterung, Tabelle E2 -- Kursangebot x Asylstatus auf die fuenf
%           Mechanismen. Kern der Erweiterung: getestet wird nicht der Ertrag
%           der Sprache, sondern der Zugang zur Behandlung selbst.
% Source  : SOEPremote jobs 244836, 244837, 244840,
%           src/stata/remote/ext_asyl_mediators_primary.do,
%           ext_asyl_mediators_secondary.do, ext_asyl_participation.do
%           Rohlistings: results/transcripts/src_stata_remote_ext_asyl_*.eml
%           Transkript: results/comparisons/ext_asyl_mediators.md
% Needs   : booktabs, tabularx (beide bereits in bachelorarbeit.sty geladen)
% Usage   : % GENERATED BY src/python/lmi/tables/ext_asyl_mediators.py
% -- DO NOT EDIT
%==============================================================================%
% ext_asyl_mediators.tex
%
% Purpose : Erweiterung, Tabelle E2 -- Kursangebot x Asylstatus auf die fuenf
%           Mechanismen. Kern der Erweiterung: getestet wird nicht der Ertrag
%           der Sprache, sondern der Zugang zur Behandlung selbst.
% Source  : SOEPremote jobs 244836, 244837, 244840,
%           src/stata/remote/ext_asyl_mediators_primary.do,
%           ext_asyl_mediators_secondary.do, ext_asyl_participation.do
%           Rohlistings: results/transcripts/src_stata_remote_ext_asyl_*.eml
%           Transkript: results/comparisons/ext_asyl_mediators.md
% Needs   : booktabs, tabularx (beide bereits in bachelorarbeit.sty geladen)
% Usage   : % GENERATED BY src/python/lmi/tables/ext_asyl_mediators.py
% -- DO NOT EDIT
%==============================================================================%
% ext_asyl_mediators.tex
%
% Purpose : Erweiterung, Tabelle E2 -- Kursangebot x Asylstatus auf die fuenf
%           Mechanismen. Kern der Erweiterung: getestet wird nicht der Ertrag
%           der Sprache, sondern der Zugang zur Behandlung selbst.
% Source  : SOEPremote jobs 244836, 244837, 244840,
%           src/stata/remote/ext_asyl_mediators_primary.do,
%           ext_asyl_mediators_secondary.do, ext_asyl_participation.do
%           Rohlistings: results/transcripts/src_stata_remote_ext_asyl_*.eml
%           Transkript: results/comparisons/ext_asyl_mediators.md
% Needs   : booktabs, tabularx (beide bereits in bachelorarbeit.sty geladen)
% Usage   : \input{tables/ext_asyl_mediators}
% Author  : Emre Suemer, 2026-09-01
%==============================================================================%

\makeatletter
\@ifundefined{NC@rewrite@Z}{%
  \newcolumntype{Z}{>{\raggedright\arraybackslash\hangindent=1.9em\hangafter=1}X}%
}{}
\makeatother

\begin{table}[htbp]
  \centering
  \caption{Marginaler Effekt des Kursangebots auf die Mechanismen,
           nach Asylstatus}
  \label{tab:ext-asyl-mechanismen}
  \footnotesize
  \setlength{\tabcolsep}{4pt}%
  \begin{tabularx}{\textwidth}{@{}Zrrrrr@{}}
    \toprule
     & Keine & & & Gemeins. & \\
    Abh\"angige Variable & Entscheidung & Anerkannt & Abgelehnt & Test $p$
     & Form \\
    \midrule
    Integrationskurs besucht & 0{,}008 (0{,}018) & 0{,}038$^{+}$ (0{,}020) & $-$0{,}035 (0{,}026) & 0{,}053 & H1 \\
    Sprachzertifikat erworben & $-$0{,}006 (0{,}016) & $-$0{,}001 (0{,}015) & $-$0{,}047$^{*}$ (0{,}019) & 0{,}105 & H1 \\
    Integrationskurs abgeschlossen & 0{,}013 (0{,}016) & 0{,}024 (0{,}021) & $-$0{,}027 (0{,}026) & 0{,}233 & H1 \\
    Kontakth\"aufigkeit zu Deutschen \textit{(Placebo)} & $-$0{,}001 (0{,}038) & $-$0{,}011 (0{,}036) & $-$0{,}065 (0{,}052) & 0{,}508 & -- \\
    Deutschkenntnisse & 0{,}048 (0{,}096) & $-$0{,}034 (0{,}099) & $-$0{,}022 (0{,}118) & 0{,}763 & -- \\
    \bottomrule
  \end{tabularx}

  \vspace{0.75ex}
  \parbox{\textwidth}{\footnotesize
    \textit{Anmerkungen:} Eigene Berechnung, IAB-BAMF-SOEP-Befragung von
    Gefl\"uchteten 2016--2019, SOEP-Core v36 (Remote Edition), gewichtet.
    
    
%    Je Modell N\,=\,5.467 Beobachtungen, 2.730 Personen, 20 Imputationen.
%    Ausgewiesen ist der durchschnittliche marginale Effekt einer Erh\"ohung
%    des Kreisangebots um eine Standardabweichung auf die jeweilige abh\"angige
%    Variable, gesch\"atzt in der Spezifikation der
%    Tabelle~\ref{tab:intermediaere-ergebnisse} zuz\"uglich der Interaktion mit
%    dem Asylstatus (\texttt{mimrgns}, Poolung nach Rubins Regeln).
    
    
%    Standardfehler in Klammern. \textbf{Rohe Einheiten, keine Prozentpunkte},
%    da die f\"unf abh\"angigen Variablen unterschiedlich skaliert sind: die
%    drei bin\"aren Variablen sind Wahrscheinlichkeiten (0{,}01\,=\,1~p.p.),
%    Deutschkenntnisse Skalenpunkte, Kontakth\"aufigkeit standardisiert.
%    Die Spalte \glqq Form\grqq{} gibt an, welcher Hypothese die Rangfolge der
%    drei Punktsch\"atzer entspricht: H1 (Motivation/Anspruch) sagt
%    abgelehnt\,$<$\,keine Entscheidung\,$<$\,anerkannt voraus, H2
%    (Rationierung) ein Maximum bei \glqq keine Entscheidung\grqq{}.
%    Kontrolliert ist die Identit\"at
%    AME(Gruppe)\,=\,$b$[Angebot]\,$+$\,$b$[Angebot\,$\times$\,Gruppe] gegen
%    \texttt{mimrgns}; gr\"o\ss{}te Abweichung zwischen beiden Rechenwegen
%    0{,}0005.
%    $^{+}$~p\,$<$\,0{,}1, $^{*}$~p\,$<$\,0{,}05, $^{**}$~p\,$<$\,0{,}01
%    (zweiseitig).
    
    
  
%  \\[0.4ex]
%  \textbf{Inferenzregel, vorab festgelegt} (Designdokument, Abschnitt~3e):
%  Prim\"ar ist je Ergebnis der gemeinsame Test beider Interaktionsterme,
%  einzelne Kontraste sind deskriptiv, dazu die Holm-Korrektur \"uber die
%  f\"unf Ergebnisse. \emph{Kein} Ergebnis erreicht p\,$<$\,0{,}05; die
%  Kursteilnahme kommt mit p\,=\,0{,}053 am n\"achsten. Die Holm-Prozedur
%  verwirft nichts: sortierte p-Werte 0{,}053; 0{,}105; 0{,}233; 0{,}508; 0{,}763 gegen die Schwellen
%  0{,}010; 0{,}013; 0{,}017; 0{,}025; 0{,}050 -- bereits der erste Vergleich scheitert, womit das Verfahren
%  endet. Berichtet werden darf, dass die Punktsch\"atzer auf der
%  Teilnahmestufe monoton mit dem rechtlichen Anspruch geordnet sind, nicht
%  aber ein einzelner signifikanter Kontrast.
%  \\[0.4ex]
%  \textbf{Zur Lesart:} Die Kursteilnahme ist der sch\"arfste Test -- ein
%  einziges kausales Glied statt vier, gemessen im Moment der Einschreibung
%  und am wenigsten imputationsabh\"angig (39 unvollst\"andige
%  Beobachtungen gegen\"uber 270 beim Kursabschluss). Das Placebo
%  verh\"alt sich wie ein Placebo: bei der Kontakth\"aufigkeit findet bereits
%  die Vorlage keinen Angebotseffekt, und eine signifikante Moderation dort
%  h\"atte nahegelegt, dass die Interaktion etwas anderes als den
%  differenzierten Kurszugang aufgreift.
    
    
    }
\end{table}

\begin{table}[htbp]
  \centering
  \caption{Kursangebot und Asylstatus: Koeffizienten der Mechanismenmodelle}
  \label{tab:ext-asyl-mechanismen-koef}
  \scriptsize
  \setlength{\tabcolsep}{3pt}%
  \begin{tabularx}{\textwidth}{@{}Zrrrrr@{}}
    \toprule
     & (1) & (2) & (3) & (4) & (5) \\
    Abh\"angige Variable & Teil- & Zerti- & Ab- & Kon- & Deutsch- \\
     & nahme & fikat & schluss & takt & kenntnisse \\
    \midrule
    \textbf{Kreisangebot an Integrationskursen, std.} & 0{,}0084 & $-$0{,}0056 & 0{,}0132 & $-$0{,}0006 & 0{,}0484 \\
     & (0{,}0181) & (0{,}0165) & (0{,}0158) & (0{,}0377) & (0{,}0964) \\[0.3ex]
    \quad $\times$ Anerkannt & 0{,}0299 & 0{,}0042 & 0{,}0112 & $-$0{,}0101 & $-$0{,}0829 \\
     & (0{,}0236) & (0{,}0190) & (0{,}0228) & (0{,}0439) & (0{,}1224) \\[0.3ex]
    \quad $\times$ Abgelehnt & $-$0{,}0431 & $-$0{,}0419$^{+}$ & $-$0{,}0402 & $-$0{,}0644 & $-$0{,}0705 \\
     & (0{,}0301) & (0{,}0234) & (0{,}0281) & (0{,}0559) & (0{,}1369) \\[0.3ex]
    Aufenthaltsdauer, in Monaten (/12) & 0{,}0361$^{+}$ & 0{,}1069$^{**}$ & 0{,}0591$^{**}$ & 0{,}1362$^{**}$ & 0{,}9839$^{**}$ \\
     & (0{,}0214) & (0{,}0161) & (0{,}0203) & (0{,}0430) & (0{,}1054) \\[0.3ex]
    Asylstatus: anerkannt & 0{,}1736$^{**}$ & 0{,}1496$^{**}$ & 0{,}1067$^{**}$ & 0{,}1252$^{+}$ & 0{,}8377$^{**}$ \\
     & (0{,}0336) & (0{,}0252) & (0{,}0310) & (0{,}0681) & (0{,}1555) \\[0.3ex]
    Asylstatus: abgelehnt & 0{,}0430 & 0{,}0773$^{**}$ & 0{,}0596$^{+}$ & 0{,}0958 & $-$0{,}1538 \\
     & (0{,}0348) & (0{,}0280) & (0{,}0332) & (0{,}0624) & (0{,}1651) \\[0.3ex]
    Konstante & 0{,}1942$^{*}$ & 0{,}0174 & 0{,}0604 & 0{,}3168$^{+}$ & 4{,}7715$^{**}$ \\
     & (0{,}0825) & (0{,}0617) & (0{,}0744) & (0{,}1793) & (0{,}4124) \\
    \midrule
    N Beobachtungen & 5.467 & 5.467 & 5.467 & 5.467 & 5.467 \\
    N Personen & 2.730 & 2.730 & 2.730 & 2.730 & 2.730 \\
    N Imputationen & 20 & 20 & 20 & 20 & 20 \\
    R\textsuperscript{2} (adj.) & 0{,}2064 & 0{,}3392 & 0{,}2030 & 0{,}1497 & 0{,}3900 \\
    \bottomrule
  \end{tabularx}

  \vspace{0.75ex}
  \parbox{\textwidth}{\footnotesize
    \textit{Anmerkungen:} Quelle und Einheiten wie
    Tabelle~\ref{tab:ext-asyl-mechanismen}. Individuelle und regionale
    Kontrollvariablen sowie fixe Effekte f\"ur Erhebungsjahr,
    Erhebungsstichprobe, Bundesland und Herkunftsland sind in jedem Modell
    enthalten und aus Platzgr\"unden nicht ausgewiesen; die vollst\"andigen
    Koeffizienten stehen in
    \texttt{build/tables/ext\_asyl\_mediators.csv}.
    Referenzkategorie des Asylstatus ist \glqq keine Entscheidung\grqq{}.
    Spalte~(1) sch\"atzt auf dem Datensatz \texttt{sample\_imputed\_part}, in
    dem \texttt{int\_course\_part} imputiert ist; die \"ubrigen Spalten auf
    \texttt{sample\_imputed}. Beide gehen aus derselben analytischen
    Stichprobe hervor, die f\"unf Spalten sind daher untereinander und mit den
    Tabellen~\ref{tab:erwerbstaetigkeit} bis~\ref{tab:mediation} vergleichbar.
    $^{+}$~p\,$<$\,0{,}1, $^{*}$~p\,$<$\,0{,}05, $^{**}$~p\,$<$\,0{,}01
    (zweiseitig).
    \\[0.4ex]
    \textbf{Die Zeile \glqq Kreisangebot\grqq{} ist nicht der Haupteffekt aus
    Tabelle~\ref{tab:intermediaere-ergebnisse}.} Mit Interaktion im Modell ist
    sie der Effekt \emph{f\"ur die Referenzgruppe} -- Personen ohne
    Entscheidung --, nicht der gepoolte Durchschnitt.}
\end{table}

% Author  : Emre Suemer, 2026-09-01
%==============================================================================%

\makeatletter
\@ifundefined{NC@rewrite@Z}{%
  \newcolumntype{Z}{>{\raggedright\arraybackslash\hangindent=1.9em\hangafter=1}X}%
}{}
\makeatother

\begin{table}[htbp]
  \centering
  \caption{Marginaler Effekt des Kursangebots auf die Mechanismen,
           nach Asylstatus}
  \label{tab:ext-asyl-mechanismen}
  \footnotesize
  \setlength{\tabcolsep}{4pt}%
  \begin{tabularx}{\textwidth}{@{}Zrrrrr@{}}
    \toprule
     & Keine & & & Gemeins. & \\
    Abh\"angige Variable & Entscheidung & Anerkannt & Abgelehnt & Test $p$
     & Form \\
    \midrule
    Integrationskurs besucht & 0{,}008 (0{,}018) & 0{,}038$^{+}$ (0{,}020) & $-$0{,}035 (0{,}026) & 0{,}053 & H1 \\
    Sprachzertifikat erworben & $-$0{,}006 (0{,}016) & $-$0{,}001 (0{,}015) & $-$0{,}047$^{*}$ (0{,}019) & 0{,}105 & H1 \\
    Integrationskurs abgeschlossen & 0{,}013 (0{,}016) & 0{,}024 (0{,}021) & $-$0{,}027 (0{,}026) & 0{,}233 & H1 \\
    Kontakth\"aufigkeit zu Deutschen \textit{(Placebo)} & $-$0{,}001 (0{,}038) & $-$0{,}011 (0{,}036) & $-$0{,}065 (0{,}052) & 0{,}508 & -- \\
    Deutschkenntnisse & 0{,}048 (0{,}096) & $-$0{,}034 (0{,}099) & $-$0{,}022 (0{,}118) & 0{,}763 & -- \\
    \bottomrule
  \end{tabularx}

  \vspace{0.75ex}
  \parbox{\textwidth}{\footnotesize
    \textit{Anmerkungen:} Eigene Berechnung, IAB-BAMF-SOEP-Befragung von
    Gefl\"uchteten 2016--2019, SOEP-Core v36 (Remote Edition), gewichtet.
    
    
%    Je Modell N\,=\,5.467 Beobachtungen, 2.730 Personen, 20 Imputationen.
%    Ausgewiesen ist der durchschnittliche marginale Effekt einer Erh\"ohung
%    des Kreisangebots um eine Standardabweichung auf die jeweilige abh\"angige
%    Variable, gesch\"atzt in der Spezifikation der
%    Tabelle~\ref{tab:intermediaere-ergebnisse} zuz\"uglich der Interaktion mit
%    dem Asylstatus (\texttt{mimrgns}, Poolung nach Rubins Regeln).
    
    
%    Standardfehler in Klammern. \textbf{Rohe Einheiten, keine Prozentpunkte},
%    da die f\"unf abh\"angigen Variablen unterschiedlich skaliert sind: die
%    drei bin\"aren Variablen sind Wahrscheinlichkeiten (0{,}01\,=\,1~p.p.),
%    Deutschkenntnisse Skalenpunkte, Kontakth\"aufigkeit standardisiert.
%    Die Spalte \glqq Form\grqq{} gibt an, welcher Hypothese die Rangfolge der
%    drei Punktsch\"atzer entspricht: H1 (Motivation/Anspruch) sagt
%    abgelehnt\,$<$\,keine Entscheidung\,$<$\,anerkannt voraus, H2
%    (Rationierung) ein Maximum bei \glqq keine Entscheidung\grqq{}.
%    Kontrolliert ist die Identit\"at
%    AME(Gruppe)\,=\,$b$[Angebot]\,$+$\,$b$[Angebot\,$\times$\,Gruppe] gegen
%    \texttt{mimrgns}; gr\"o\ss{}te Abweichung zwischen beiden Rechenwegen
%    0{,}0005.
%    $^{+}$~p\,$<$\,0{,}1, $^{*}$~p\,$<$\,0{,}05, $^{**}$~p\,$<$\,0{,}01
%    (zweiseitig).
    
    
  
%  \\[0.4ex]
%  \textbf{Inferenzregel, vorab festgelegt} (Designdokument, Abschnitt~3e):
%  Prim\"ar ist je Ergebnis der gemeinsame Test beider Interaktionsterme,
%  einzelne Kontraste sind deskriptiv, dazu die Holm-Korrektur \"uber die
%  f\"unf Ergebnisse. \emph{Kein} Ergebnis erreicht p\,$<$\,0{,}05; die
%  Kursteilnahme kommt mit p\,=\,0{,}053 am n\"achsten. Die Holm-Prozedur
%  verwirft nichts: sortierte p-Werte 0{,}053; 0{,}105; 0{,}233; 0{,}508; 0{,}763 gegen die Schwellen
%  0{,}010; 0{,}013; 0{,}017; 0{,}025; 0{,}050 -- bereits der erste Vergleich scheitert, womit das Verfahren
%  endet. Berichtet werden darf, dass die Punktsch\"atzer auf der
%  Teilnahmestufe monoton mit dem rechtlichen Anspruch geordnet sind, nicht
%  aber ein einzelner signifikanter Kontrast.
%  \\[0.4ex]
%  \textbf{Zur Lesart:} Die Kursteilnahme ist der sch\"arfste Test -- ein
%  einziges kausales Glied statt vier, gemessen im Moment der Einschreibung
%  und am wenigsten imputationsabh\"angig (39 unvollst\"andige
%  Beobachtungen gegen\"uber 270 beim Kursabschluss). Das Placebo
%  verh\"alt sich wie ein Placebo: bei der Kontakth\"aufigkeit findet bereits
%  die Vorlage keinen Angebotseffekt, und eine signifikante Moderation dort
%  h\"atte nahegelegt, dass die Interaktion etwas anderes als den
%  differenzierten Kurszugang aufgreift.
    
    
    }
\end{table}

\begin{table}[htbp]
  \centering
  \caption{Kursangebot und Asylstatus: Koeffizienten der Mechanismenmodelle}
  \label{tab:ext-asyl-mechanismen-koef}
  \scriptsize
  \setlength{\tabcolsep}{3pt}%
  \begin{tabularx}{\textwidth}{@{}Zrrrrr@{}}
    \toprule
     & (1) & (2) & (3) & (4) & (5) \\
    Abh\"angige Variable & Teil- & Zerti- & Ab- & Kon- & Deutsch- \\
     & nahme & fikat & schluss & takt & kenntnisse \\
    \midrule
    \textbf{Kreisangebot an Integrationskursen, std.} & 0{,}0084 & $-$0{,}0056 & 0{,}0132 & $-$0{,}0006 & 0{,}0484 \\
     & (0{,}0181) & (0{,}0165) & (0{,}0158) & (0{,}0377) & (0{,}0964) \\[0.3ex]
    \quad $\times$ Anerkannt & 0{,}0299 & 0{,}0042 & 0{,}0112 & $-$0{,}0101 & $-$0{,}0829 \\
     & (0{,}0236) & (0{,}0190) & (0{,}0228) & (0{,}0439) & (0{,}1224) \\[0.3ex]
    \quad $\times$ Abgelehnt & $-$0{,}0431 & $-$0{,}0419$^{+}$ & $-$0{,}0402 & $-$0{,}0644 & $-$0{,}0705 \\
     & (0{,}0301) & (0{,}0234) & (0{,}0281) & (0{,}0559) & (0{,}1369) \\[0.3ex]
    Aufenthaltsdauer, in Monaten (/12) & 0{,}0361$^{+}$ & 0{,}1069$^{**}$ & 0{,}0591$^{**}$ & 0{,}1362$^{**}$ & 0{,}9839$^{**}$ \\
     & (0{,}0214) & (0{,}0161) & (0{,}0203) & (0{,}0430) & (0{,}1054) \\[0.3ex]
    Asylstatus: anerkannt & 0{,}1736$^{**}$ & 0{,}1496$^{**}$ & 0{,}1067$^{**}$ & 0{,}1252$^{+}$ & 0{,}8377$^{**}$ \\
     & (0{,}0336) & (0{,}0252) & (0{,}0310) & (0{,}0681) & (0{,}1555) \\[0.3ex]
    Asylstatus: abgelehnt & 0{,}0430 & 0{,}0773$^{**}$ & 0{,}0596$^{+}$ & 0{,}0958 & $-$0{,}1538 \\
     & (0{,}0348) & (0{,}0280) & (0{,}0332) & (0{,}0624) & (0{,}1651) \\[0.3ex]
    Konstante & 0{,}1942$^{*}$ & 0{,}0174 & 0{,}0604 & 0{,}3168$^{+}$ & 4{,}7715$^{**}$ \\
     & (0{,}0825) & (0{,}0617) & (0{,}0744) & (0{,}1793) & (0{,}4124) \\
    \midrule
    N Beobachtungen & 5.467 & 5.467 & 5.467 & 5.467 & 5.467 \\
    N Personen & 2.730 & 2.730 & 2.730 & 2.730 & 2.730 \\
    N Imputationen & 20 & 20 & 20 & 20 & 20 \\
    R\textsuperscript{2} (adj.) & 0{,}2064 & 0{,}3392 & 0{,}2030 & 0{,}1497 & 0{,}3900 \\
    \bottomrule
  \end{tabularx}

  \vspace{0.75ex}
  \parbox{\textwidth}{\footnotesize
    \textit{Anmerkungen:} Quelle und Einheiten wie
    Tabelle~\ref{tab:ext-asyl-mechanismen}. Individuelle und regionale
    Kontrollvariablen sowie fixe Effekte f\"ur Erhebungsjahr,
    Erhebungsstichprobe, Bundesland und Herkunftsland sind in jedem Modell
    enthalten und aus Platzgr\"unden nicht ausgewiesen; die vollst\"andigen
    Koeffizienten stehen in
    \texttt{build/tables/ext\_asyl\_mediators.csv}.
    Referenzkategorie des Asylstatus ist \glqq keine Entscheidung\grqq{}.
    Spalte~(1) sch\"atzt auf dem Datensatz \texttt{sample\_imputed\_part}, in
    dem \texttt{int\_course\_part} imputiert ist; die \"ubrigen Spalten auf
    \texttt{sample\_imputed}. Beide gehen aus derselben analytischen
    Stichprobe hervor, die f\"unf Spalten sind daher untereinander und mit den
    Tabellen~\ref{tab:erwerbstaetigkeit} bis~\ref{tab:mediation} vergleichbar.
    $^{+}$~p\,$<$\,0{,}1, $^{*}$~p\,$<$\,0{,}05, $^{**}$~p\,$<$\,0{,}01
    (zweiseitig).
    \\[0.4ex]
    \textbf{Die Zeile \glqq Kreisangebot\grqq{} ist nicht der Haupteffekt aus
    Tabelle~\ref{tab:intermediaere-ergebnisse}.} Mit Interaktion im Modell ist
    sie der Effekt \emph{f\"ur die Referenzgruppe} -- Personen ohne
    Entscheidung --, nicht der gepoolte Durchschnitt.}
\end{table}

% Author  : Emre Suemer, 2026-09-01
%==============================================================================%

\makeatletter
\@ifundefined{NC@rewrite@Z}{%
  \newcolumntype{Z}{>{\raggedright\arraybackslash\hangindent=1.9em\hangafter=1}X}%
}{}
\makeatother

\begin{table}[htbp]
  \centering
  \caption{Marginaler Effekt des Kursangebots auf die Mechanismen,
           nach Asylstatus}
  \label{tab:ext-asyl-mechanismen}
  \footnotesize
  \setlength{\tabcolsep}{4pt}%
  \begin{tabularx}{\textwidth}{@{}Zrrrrr@{}}
    \toprule
     & Keine & & & Gemeins. & \\
    Abh\"angige Variable & Entscheidung & Anerkannt & Abgelehnt & Test $p$
     & Form \\
    \midrule
    Integrationskurs besucht & 0{,}008 (0{,}018) & 0{,}038$^{+}$ (0{,}020) & $-$0{,}035 (0{,}026) & 0{,}053 & H1 \\
    Sprachzertifikat erworben & $-$0{,}006 (0{,}016) & $-$0{,}001 (0{,}015) & $-$0{,}047$^{*}$ (0{,}019) & 0{,}105 & H1 \\
    Integrationskurs abgeschlossen & 0{,}013 (0{,}016) & 0{,}024 (0{,}021) & $-$0{,}027 (0{,}026) & 0{,}233 & H1 \\
    Kontakth\"aufigkeit zu Deutschen \textit{(Placebo)} & $-$0{,}001 (0{,}038) & $-$0{,}011 (0{,}036) & $-$0{,}065 (0{,}052) & 0{,}508 & -- \\
    Deutschkenntnisse & 0{,}048 (0{,}096) & $-$0{,}034 (0{,}099) & $-$0{,}022 (0{,}118) & 0{,}763 & -- \\
    \bottomrule
  \end{tabularx}

  \vspace{0.75ex}
  \parbox{\textwidth}{\footnotesize
    \textit{Anmerkungen:} Eigene Berechnung, IAB-BAMF-SOEP-Befragung von
    Gefl\"uchteten 2016--2019, SOEP-Core v36 (Remote Edition), gewichtet.
    
    
%    Je Modell N\,=\,5.467 Beobachtungen, 2.730 Personen, 20 Imputationen.
%    Ausgewiesen ist der durchschnittliche marginale Effekt einer Erh\"ohung
%    des Kreisangebots um eine Standardabweichung auf die jeweilige abh\"angige
%    Variable, gesch\"atzt in der Spezifikation der
%    Tabelle~\ref{tab:intermediaere-ergebnisse} zuz\"uglich der Interaktion mit
%    dem Asylstatus (\texttt{mimrgns}, Poolung nach Rubins Regeln).
    
    
%    Standardfehler in Klammern. \textbf{Rohe Einheiten, keine Prozentpunkte},
%    da die f\"unf abh\"angigen Variablen unterschiedlich skaliert sind: die
%    drei bin\"aren Variablen sind Wahrscheinlichkeiten (0{,}01\,=\,1~p.p.),
%    Deutschkenntnisse Skalenpunkte, Kontakth\"aufigkeit standardisiert.
%    Die Spalte \glqq Form\grqq{} gibt an, welcher Hypothese die Rangfolge der
%    drei Punktsch\"atzer entspricht: H1 (Motivation/Anspruch) sagt
%    abgelehnt\,$<$\,keine Entscheidung\,$<$\,anerkannt voraus, H2
%    (Rationierung) ein Maximum bei \glqq keine Entscheidung\grqq{}.
%    Kontrolliert ist die Identit\"at
%    AME(Gruppe)\,=\,$b$[Angebot]\,$+$\,$b$[Angebot\,$\times$\,Gruppe] gegen
%    \texttt{mimrgns}; gr\"o\ss{}te Abweichung zwischen beiden Rechenwegen
%    0{,}0005.
%    $^{+}$~p\,$<$\,0{,}1, $^{*}$~p\,$<$\,0{,}05, $^{**}$~p\,$<$\,0{,}01
%    (zweiseitig).
    
    
  
%  \\[0.4ex]
%  \textbf{Inferenzregel, vorab festgelegt} (Designdokument, Abschnitt~3e):
%  Prim\"ar ist je Ergebnis der gemeinsame Test beider Interaktionsterme,
%  einzelne Kontraste sind deskriptiv, dazu die Holm-Korrektur \"uber die
%  f\"unf Ergebnisse. \emph{Kein} Ergebnis erreicht p\,$<$\,0{,}05; die
%  Kursteilnahme kommt mit p\,=\,0{,}053 am n\"achsten. Die Holm-Prozedur
%  verwirft nichts: sortierte p-Werte 0{,}053; 0{,}105; 0{,}233; 0{,}508; 0{,}763 gegen die Schwellen
%  0{,}010; 0{,}013; 0{,}017; 0{,}025; 0{,}050 -- bereits der erste Vergleich scheitert, womit das Verfahren
%  endet. Berichtet werden darf, dass die Punktsch\"atzer auf der
%  Teilnahmestufe monoton mit dem rechtlichen Anspruch geordnet sind, nicht
%  aber ein einzelner signifikanter Kontrast.
%  \\[0.4ex]
%  \textbf{Zur Lesart:} Die Kursteilnahme ist der sch\"arfste Test -- ein
%  einziges kausales Glied statt vier, gemessen im Moment der Einschreibung
%  und am wenigsten imputationsabh\"angig (39 unvollst\"andige
%  Beobachtungen gegen\"uber 270 beim Kursabschluss). Das Placebo
%  verh\"alt sich wie ein Placebo: bei der Kontakth\"aufigkeit findet bereits
%  die Vorlage keinen Angebotseffekt, und eine signifikante Moderation dort
%  h\"atte nahegelegt, dass die Interaktion etwas anderes als den
%  differenzierten Kurszugang aufgreift.
    
    
    }
\end{table}

\begin{table}[htbp]
  \centering
  \caption{Kursangebot und Asylstatus: Koeffizienten der Mechanismenmodelle}
  \label{tab:ext-asyl-mechanismen-koef}
  \scriptsize
  \setlength{\tabcolsep}{3pt}%
  \begin{tabularx}{\textwidth}{@{}Zrrrrr@{}}
    \toprule
     & (1) & (2) & (3) & (4) & (5) \\
    Abh\"angige Variable & Teil- & Zerti- & Ab- & Kon- & Deutsch- \\
     & nahme & fikat & schluss & takt & kenntnisse \\
    \midrule
    \textbf{Kreisangebot an Integrationskursen, std.} & 0{,}0084 & $-$0{,}0056 & 0{,}0132 & $-$0{,}0006 & 0{,}0484 \\
     & (0{,}0181) & (0{,}0165) & (0{,}0158) & (0{,}0377) & (0{,}0964) \\[0.3ex]
    \quad $\times$ Anerkannt & 0{,}0299 & 0{,}0042 & 0{,}0112 & $-$0{,}0101 & $-$0{,}0829 \\
     & (0{,}0236) & (0{,}0190) & (0{,}0228) & (0{,}0439) & (0{,}1224) \\[0.3ex]
    \quad $\times$ Abgelehnt & $-$0{,}0431 & $-$0{,}0419$^{+}$ & $-$0{,}0402 & $-$0{,}0644 & $-$0{,}0705 \\
     & (0{,}0301) & (0{,}0234) & (0{,}0281) & (0{,}0559) & (0{,}1369) \\[0.3ex]
    Aufenthaltsdauer, in Monaten (/12) & 0{,}0361$^{+}$ & 0{,}1069$^{**}$ & 0{,}0591$^{**}$ & 0{,}1362$^{**}$ & 0{,}9839$^{**}$ \\
     & (0{,}0214) & (0{,}0161) & (0{,}0203) & (0{,}0430) & (0{,}1054) \\[0.3ex]
    Asylstatus: anerkannt & 0{,}1736$^{**}$ & 0{,}1496$^{**}$ & 0{,}1067$^{**}$ & 0{,}1252$^{+}$ & 0{,}8377$^{**}$ \\
     & (0{,}0336) & (0{,}0252) & (0{,}0310) & (0{,}0681) & (0{,}1555) \\[0.3ex]
    Asylstatus: abgelehnt & 0{,}0430 & 0{,}0773$^{**}$ & 0{,}0596$^{+}$ & 0{,}0958 & $-$0{,}1538 \\
     & (0{,}0348) & (0{,}0280) & (0{,}0332) & (0{,}0624) & (0{,}1651) \\[0.3ex]
    Konstante & 0{,}1942$^{*}$ & 0{,}0174 & 0{,}0604 & 0{,}3168$^{+}$ & 4{,}7715$^{**}$ \\
     & (0{,}0825) & (0{,}0617) & (0{,}0744) & (0{,}1793) & (0{,}4124) \\
    \midrule
    N Beobachtungen & 5.467 & 5.467 & 5.467 & 5.467 & 5.467 \\
    N Personen & 2.730 & 2.730 & 2.730 & 2.730 & 2.730 \\
    N Imputationen & 20 & 20 & 20 & 20 & 20 \\
    R\textsuperscript{2} (adj.) & 0{,}2064 & 0{,}3392 & 0{,}2030 & 0{,}1497 & 0{,}3900 \\
    \bottomrule
  \end{tabularx}

  \vspace{0.75ex}
  \parbox{\textwidth}{\footnotesize
    \textit{Anmerkungen:} Quelle und Einheiten wie
    Tabelle~\ref{tab:ext-asyl-mechanismen}. Individuelle und regionale
    Kontrollvariablen sowie fixe Effekte f\"ur Erhebungsjahr,
    Erhebungsstichprobe, Bundesland und Herkunftsland sind in jedem Modell
    enthalten und aus Platzgr\"unden nicht ausgewiesen; die vollst\"andigen
    Koeffizienten stehen in
    \texttt{build/tables/ext\_asyl\_mediators.csv}.
    Referenzkategorie des Asylstatus ist \glqq keine Entscheidung\grqq{}.
    Spalte~(1) sch\"atzt auf dem Datensatz \texttt{sample\_imputed\_part}, in
    dem \texttt{int\_course\_part} imputiert ist; die \"ubrigen Spalten auf
    \texttt{sample\_imputed}. Beide gehen aus derselben analytischen
    Stichprobe hervor, die f\"unf Spalten sind daher untereinander und mit den
    Tabellen~\ref{tab:erwerbstaetigkeit} bis~\ref{tab:mediation} vergleichbar.
    $^{+}$~p\,$<$\,0{,}1, $^{*}$~p\,$<$\,0{,}05, $^{**}$~p\,$<$\,0{,}01
    (zweiseitig).
    \\[0.4ex]
    \textbf{Die Zeile \glqq Kreisangebot\grqq{} ist nicht der Haupteffekt aus
    Tabelle~\ref{tab:intermediaere-ergebnisse}.} Mit Interaktion im Modell ist
    sie der Effekt \emph{f\"ur die Referenzgruppe} -- Personen ohne
    Entscheidung --, nicht der gepoolte Durchschnitt.}
\end{table}

% Author  : Emre Suemer, 2026-09-01
%==============================================================================%

\makeatletter
\@ifundefined{NC@rewrite@Z}{%
  \newcolumntype{Z}{>{\raggedright\arraybackslash\hangindent=1.9em\hangafter=1}X}%
}{}
\makeatother

\begin{table}[htbp]
  \centering
  \caption{Marginaler Effekt des Kursangebots auf die Mechanismen,
           nach Asylstatus}
  \label{tab:ext-asyl-mechanismen}
  \footnotesize
  \setlength{\tabcolsep}{4pt}%
  \begin{tabularx}{\textwidth}{@{}Zrrrrr@{}}
    \toprule
     & Keine & & & Gemeins. & \\
    Abh\"angige Variable & Entscheidung & Anerkannt & Abgelehnt & Test $p$
     & Form \\
    \midrule
    Integrationskurs besucht & 0{,}008 (0{,}018) & 0{,}038$^{+}$ (0{,}020) & $-$0{,}035 (0{,}026) & 0{,}053 & H1 \\
    Sprachzertifikat erworben & $-$0{,}006 (0{,}016) & $-$0{,}001 (0{,}015) & $-$0{,}047$^{*}$ (0{,}019) & 0{,}105 & H1 \\
    Integrationskurs abgeschlossen & 0{,}013 (0{,}016) & 0{,}024 (0{,}021) & $-$0{,}027 (0{,}026) & 0{,}233 & H1 \\
    Kontakth\"aufigkeit zu Deutschen \textit{(Placebo)} & $-$0{,}001 (0{,}038) & $-$0{,}011 (0{,}036) & $-$0{,}065 (0{,}052) & 0{,}508 & -- \\
    Deutschkenntnisse & 0{,}048 (0{,}096) & $-$0{,}034 (0{,}099) & $-$0{,}022 (0{,}118) & 0{,}763 & -- \\
    \bottomrule
  \end{tabularx}

  \vspace{0.75ex}
  \parbox{\textwidth}{\footnotesize
    \textit{Anmerkungen:} Eigene Berechnung, IAB-BAMF-SOEP-Befragung von
    Gefl\"uchteten 2016--2019, SOEP-Core v36 (Remote Edition), gewichtet.
    
    
%    Je Modell N\,=\,5.467 Beobachtungen, 2.730 Personen, 20 Imputationen.
%    Ausgewiesen ist der durchschnittliche marginale Effekt einer Erh\"ohung
%    des Kreisangebots um eine Standardabweichung auf die jeweilige abh\"angige
%    Variable, gesch\"atzt in der Spezifikation der
%    Tabelle~\ref{tab:intermediaere-ergebnisse} zuz\"uglich der Interaktion mit
%    dem Asylstatus (\texttt{mimrgns}, Poolung nach Rubins Regeln).
    
    
%    Standardfehler in Klammern. \textbf{Rohe Einheiten, keine Prozentpunkte},
%    da die f\"unf abh\"angigen Variablen unterschiedlich skaliert sind: die
%    drei bin\"aren Variablen sind Wahrscheinlichkeiten (0{,}01\,=\,1~p.p.),
%    Deutschkenntnisse Skalenpunkte, Kontakth\"aufigkeit standardisiert.
%    Die Spalte \glqq Form\grqq{} gibt an, welcher Hypothese die Rangfolge der
%    drei Punktsch\"atzer entspricht: H1 (Motivation/Anspruch) sagt
%    abgelehnt\,$<$\,keine Entscheidung\,$<$\,anerkannt voraus, H2
%    (Rationierung) ein Maximum bei \glqq keine Entscheidung\grqq{}.
%    Kontrolliert ist die Identit\"at
%    AME(Gruppe)\,=\,$b$[Angebot]\,$+$\,$b$[Angebot\,$\times$\,Gruppe] gegen
%    \texttt{mimrgns}; gr\"o\ss{}te Abweichung zwischen beiden Rechenwegen
%    0{,}0005.
%    $^{+}$~p\,$<$\,0{,}1, $^{*}$~p\,$<$\,0{,}05, $^{**}$~p\,$<$\,0{,}01
%    (zweiseitig).
    
    
  
%  \\[0.4ex]
%  \textbf{Inferenzregel, vorab festgelegt} (Designdokument, Abschnitt~3e):
%  Prim\"ar ist je Ergebnis der gemeinsame Test beider Interaktionsterme,
%  einzelne Kontraste sind deskriptiv, dazu die Holm-Korrektur \"uber die
%  f\"unf Ergebnisse. \emph{Kein} Ergebnis erreicht p\,$<$\,0{,}05; die
%  Kursteilnahme kommt mit p\,=\,0{,}053 am n\"achsten. Die Holm-Prozedur
%  verwirft nichts: sortierte p-Werte 0{,}053; 0{,}105; 0{,}233; 0{,}508; 0{,}763 gegen die Schwellen
%  0{,}010; 0{,}013; 0{,}017; 0{,}025; 0{,}050 -- bereits der erste Vergleich scheitert, womit das Verfahren
%  endet. Berichtet werden darf, dass die Punktsch\"atzer auf der
%  Teilnahmestufe monoton mit dem rechtlichen Anspruch geordnet sind, nicht
%  aber ein einzelner signifikanter Kontrast.
%  \\[0.4ex]
%  \textbf{Zur Lesart:} Die Kursteilnahme ist der sch\"arfste Test -- ein
%  einziges kausales Glied statt vier, gemessen im Moment der Einschreibung
%  und am wenigsten imputationsabh\"angig (39 unvollst\"andige
%  Beobachtungen gegen\"uber 270 beim Kursabschluss). Das Placebo
%  verh\"alt sich wie ein Placebo: bei der Kontakth\"aufigkeit findet bereits
%  die Vorlage keinen Angebotseffekt, und eine signifikante Moderation dort
%  h\"atte nahegelegt, dass die Interaktion etwas anderes als den
%  differenzierten Kurszugang aufgreift.
    
    
    }
\end{table}

\begin{table}[htbp]
  \centering
  \caption{Kursangebot und Asylstatus: Koeffizienten der Mechanismenmodelle}
  \label{tab:ext-asyl-mechanismen-koef}
  \scriptsize
  \setlength{\tabcolsep}{3pt}%
  \begin{tabularx}{\textwidth}{@{}Zrrrrr@{}}
    \toprule
     & (1) & (2) & (3) & (4) & (5) \\
    Abh\"angige Variable & Teil- & Zerti- & Ab- & Kon- & Deutsch- \\
     & nahme & fikat & schluss & takt & kenntnisse \\
    \midrule
    \textbf{Kreisangebot an Integrationskursen, std.} & 0{,}0084 & $-$0{,}0056 & 0{,}0132 & $-$0{,}0006 & 0{,}0484 \\
     & (0{,}0181) & (0{,}0165) & (0{,}0158) & (0{,}0377) & (0{,}0964) \\[0.3ex]
    \quad $\times$ Anerkannt & 0{,}0299 & 0{,}0042 & 0{,}0112 & $-$0{,}0101 & $-$0{,}0829 \\
     & (0{,}0236) & (0{,}0190) & (0{,}0228) & (0{,}0439) & (0{,}1224) \\[0.3ex]
    \quad $\times$ Abgelehnt & $-$0{,}0431 & $-$0{,}0419$^{+}$ & $-$0{,}0402 & $-$0{,}0644 & $-$0{,}0705 \\
     & (0{,}0301) & (0{,}0234) & (0{,}0281) & (0{,}0559) & (0{,}1369) \\[0.3ex]
    Aufenthaltsdauer, in Monaten (/12) & 0{,}0361$^{+}$ & 0{,}1069$^{**}$ & 0{,}0591$^{**}$ & 0{,}1362$^{**}$ & 0{,}9839$^{**}$ \\
     & (0{,}0214) & (0{,}0161) & (0{,}0203) & (0{,}0430) & (0{,}1054) \\[0.3ex]
    Asylstatus: anerkannt & 0{,}1736$^{**}$ & 0{,}1496$^{**}$ & 0{,}1067$^{**}$ & 0{,}1252$^{+}$ & 0{,}8377$^{**}$ \\
     & (0{,}0336) & (0{,}0252) & (0{,}0310) & (0{,}0681) & (0{,}1555) \\[0.3ex]
    Asylstatus: abgelehnt & 0{,}0430 & 0{,}0773$^{**}$ & 0{,}0596$^{+}$ & 0{,}0958 & $-$0{,}1538 \\
     & (0{,}0348) & (0{,}0280) & (0{,}0332) & (0{,}0624) & (0{,}1651) \\[0.3ex]
    Konstante & 0{,}1942$^{*}$ & 0{,}0174 & 0{,}0604 & 0{,}3168$^{+}$ & 4{,}7715$^{**}$ \\
     & (0{,}0825) & (0{,}0617) & (0{,}0744) & (0{,}1793) & (0{,}4124) \\
    \midrule
    N Beobachtungen & 5.467 & 5.467 & 5.467 & 5.467 & 5.467 \\
    N Personen & 2.730 & 2.730 & 2.730 & 2.730 & 2.730 \\
    N Imputationen & 20 & 20 & 20 & 20 & 20 \\
    R\textsuperscript{2} (adj.) & 0{,}2064 & 0{,}3392 & 0{,}2030 & 0{,}1497 & 0{,}3900 \\
    \bottomrule
  \end{tabularx}

  \vspace{0.75ex}
  \parbox{\textwidth}{\footnotesize
    \textit{Anmerkungen:} Quelle und Einheiten wie
    Tabelle~\ref{tab:ext-asyl-mechanismen}. Individuelle und regionale
    Kontrollvariablen sowie fixe Effekte f\"ur Erhebungsjahr,
    Erhebungsstichprobe, Bundesland und Herkunftsland sind in jedem Modell
    enthalten und aus Platzgr\"unden nicht ausgewiesen; die vollst\"andigen
    Koeffizienten stehen in
    \texttt{build/tables/ext\_asyl\_mediators.csv}.
    Referenzkategorie des Asylstatus ist \glqq keine Entscheidung\grqq{}.
    Spalte~(1) sch\"atzt auf dem Datensatz \texttt{sample\_imputed\_part}, in
    dem \texttt{int\_course\_part} imputiert ist; die \"ubrigen Spalten auf
    \texttt{sample\_imputed}. Beide gehen aus derselben analytischen
    Stichprobe hervor, die f\"unf Spalten sind daher untereinander und mit den
    Tabellen~\ref{tab:erwerbstaetigkeit} bis~\ref{tab:mediation} vergleichbar.
    $^{+}$~p\,$<$\,0{,}1, $^{*}$~p\,$<$\,0{,}05, $^{**}$~p\,$<$\,0{,}01
    (zweiseitig).
    \\[0.4ex]
    \textbf{Die Zeile \glqq Kreisangebot\grqq{} ist nicht der Haupteffekt aus
    Tabelle~\ref{tab:intermediaere-ergebnisse}.} Mit Interaktion im Modell ist
    sie der Effekt \emph{f\"ur die Referenzgruppe} -- Personen ohne
    Entscheidung --, nicht der gepoolte Durchschnitt.}
\end{table}
